\documentclass[12pt]{article}
\usepackage[utf8]{inputenc}
\usepackage[T1]{fontenc}
\usepackage{amsmath}
\usepackage{amsfonts}
\usepackage{amssymb}
\usepackage{hyperref}
\usepackage{geometry}
\geometry{
 a4paper,
 margin=1in,
}
\title{\textbf{Lab ISO Documentation (Ubuntu 22.04 LTS)}}
\author{Samnun Azfar, Jr. Lecturer, CSE \\ Samia Nawsheen, Jr. Lecturer, CSE}
\date{November 2026}

\begin{document}

\maketitle

\section{User Accounts}
Two user accounts are pre-configured on the system.

\subsection*{Standard User (For Students)}
\begin{itemize}
    \item \textbf{Username:} \texttt{cse}
    \item \textbf{Password:} \texttt{123}
    \item \textbf{Notes:} This is the primary user account for student's lab work. This user has full \texttt{sudo} privileges for lab convenience.
\end{itemize}

\subsection*{Administrator Account (For Faculties)}
\begin{itemize}
    \item \textbf{Username:} \texttt{admin}
    \item \textbf{Password:} \texttt{1ut01ccs3}
    \item \textbf{Notes:} This account has full \texttt{sudo} privileges and is the only account permitted to log in via SSH for remote administration. \textbf{This password must be kept confidential and should not be shared with students.}
\end{itemize}

\section{Lab Monitoring (Epoptes), Clean-up}
\textbf{Epoptes} is a classroom monitoring software.
\begin{itemize}
    \item \textbf{Connection Script:} There is a script named \texttt{epoptes\_connect.sh} in the directory \texttt{\string~/Documents/Scripts}. Running that script will update all the turned-on PCs under the current subnetwork and connect them to the teacher PC.
    \item \textbf{Usage:} Afterwards, all PCs can be monitored (or taken access of) using the \texttt{epoptes} Program. 
    \item \textbf{Clean-up Script:} The \texttt{\string~/Documents/Scripts} directory will also contain a script named \texttt{clear\_files.sh} which will clear all the non-essential (e.g., task files) of all PCs currently turned on and on the same subnet.
    \item \textbf{Custom Configuration:} Any other SSH-based configurations can be set-up (e.g., installing specific packages, restricting certain actions) on demand.
\end{itemize}

\section{Pre-installed Software \& Configuration}

\subsection{Core Development Suite}
\begin{itemize}
    \item \textbf{Browsers:} Firefox, Google Chrome
    \item \textbf{Compilers \& Runtimes:} \texttt{gcc}, \texttt{g++}, \texttt{python2}, \texttt{python3}, \texttt{java}
    \item \textbf{Lightweight Editors:} Sublime Text, Visual Studio Code, NeoVim, Geany
    \item \textbf{Full IDEs:} Code::Blocks, PyCharm, Eclipse, NetBeans, IntelliJ
\end{itemize}

\subsection{Dev Tools}
\begin{itemize}
    \item \textbf{Git:} Installed and available.
    \item \textbf{Conda:} A system-wide installation of Conda is present. The base environment is set to activate automatically by default in new terminal sessions.
    \item \textbf{JupyterLab:} Can be launched from the applications menu (via a desktop launcher) or by running \texttt{jupyter lab} in the terminal.
\end{itemize}

\subsection{Databases \& GUIs}
\begin{itemize}
    \item \textbf{Services:} Both \texttt{postgresql} and \texttt{mysql} services are enabled and start automatically on boot.
    \item \textbf{Access:} They are configured for \textbf{passwordless local access}. Users can connect by simply typing \texttt{psql} or \texttt{mysql} in the terminal.
    \item \textbf{GUI Tools:} \texttt{pgAdmin} and \texttt{MySQL Workbench} are installed. No default connections are pre-configured.
\end{itemize}

\subsection{Virtualization}
\begin{itemize}
    \item \textbf{Docker:} The Docker service starts automatically on boot. The \texttt{cse} user can run all docker commands using \texttt{sudo}.
    \item \textbf{Virtual-Machine Manager:} Virtual-Machine Manager (QEMU/KVM) is installed, allowing for the creation and management of virtual machines.
\end{itemize}

\subsection{Additional Software}
\begin{itemize}
    \item Blender
    \item CISCO Packet Tracer (ran by the command \texttt{packettracer} on the terminal)
    \item Draw.io
\end{itemize}

\subsection{Utilities}
\begin{itemize}
    \item \textbf{Gpaste (Clipboard Manager):} The Gpaste clipboard manager is installed and active. This utility logs clipboard history, which can be accessed to help assist or catch potential offenses.
\end{itemize}

\end{document}